% ---- theorem environments (thmtools, so that cleveref names each kind correctly)
% main_pt.tex (the Portuguese translation) defines \versaopt before reading this file.
\ifdefined\versaopt
\declaretheorem[name=Teorema,refname={Teorema,Teoremas},numberwithin=section]{theorem}
\declaretheorem[name=Proposição,refname={Proposição,Proposições},sibling=theorem]{proposition}
\declaretheorem[name=Lema,refname={Lema,Lemas},sibling=theorem]{lemma}
\declaretheorem[name=Corolário,refname={Corolário,Corolários},sibling=theorem]{corollary}
\declaretheorem[name=Definição,refname={Definição,Definições},sibling=theorem,style=definition]{definition}
\declaretheorem[name=Observação,refname={Observação,Observações},sibling=theorem,style=remark]{remark}
\declaretheorem[name=Teorema Principal,numbered=no]{mainthm}
\declaretheorem[name=Teorema,numbered=no]{introthm}
\else
\declaretheorem[name=Theorem,numberwithin=section]{theorem}
\declaretheorem[name=Proposition,sibling=theorem]{proposition}
\declaretheorem[name=Lemma,sibling=theorem]{lemma}
\declaretheorem[name=Corollary,sibling=theorem]{corollary}
\declaretheorem[name=Definition,sibling=theorem,style=definition]{definition}
\declaretheorem[name=Remark,sibling=theorem,style=remark]{remark}
\declaretheorem[name=Main Theorem,numbered=no]{mainthm}
\declaretheorem[name=Theorem,numbered=no]{introthm}
\fi

% ---- drafting aids (remove before submission)
\newcommand{\todo}[1]{{\color{BrickRed}[\textbf{TODO:} #1]}}
\newcommand{\src}[1]{{\color{Gray}\footnotesize[\texttt{#1}]}}

% ---- notation
\newcommand{\R}{\mathbb{R}}
\newcommand{\C}{\mathbb{C}}
\newcommand{\Z}{\mathbb{Z}}
\renewcommand{\Re}{\operatorname{Re}}
\renewcommand{\Im}{\operatorname{Im}}
\newcommand{\gbar}{\bar g}
\newcommand{\phibar}{\bar\phi}
\newcommand{\omegabar}{\bar\omega}
\newcommand{\Mhat}{\widehat M}
\newcommand{\Lop}{L}             % linearized RT family
\newcommand{\Cop}{C}             % linearized constraints
\newcommand{\Kop}{K}             % constraint propagation operator
\newcommand{\Lie}{\mathcal{L}}
\newcommand{\sphys}{s_{\ast}}    % the physical root
\newcommand{\sK}{s_{K}}
\newcommand{\sB}{s_{B}}
\newcommand{\Dom}{\mathcal{D}}
\newcommand{\Yp}{Y_{+}}
\newcommand{\mult}{\operatorname{m}}

% DOI links in the bibliography (amsplain prints the note field)
\newcommand{\doi}[1]{\href{https://doi.org/#1}{\texttt{doi:#1}}}
